\documentclass[11pt,a4paper]{article}
\usepackage[T1]{fontenc}\usepackage{amsmath,amssymb}\usepackage[margin=2.5cm]{geometry}
\title{Creeping Flow in a Lid-Driven Square Cavity}\author{G. Fluid and H. Flow}\date{}
\begin{document}\maketitle
\section*{1. Problem description}
We consider the steady flow of an incompressible Newtonian fluid in the unit square cavity $\Omega=(0,1)\times(0,1)$. The upper wall, $y=1$, slides in its own plane with a constant velocity, whereas the other three walls are at rest. The Reynolds number is so small that inertia is negligible.

\section*{2. Governing equations}
The velocity $u=(u_x,u_y)$ and the pressure $p$ satisfy the Stokes equations

\begin{equation}
-\nabla\cdot\left(2\nu\,\varepsilon(u)\right) + \nabla p = 0 \quad \text{in } \Omega,
\end{equation}
\begin{equation}
\nabla\cdot u = 0 \quad \text{in } \Omega,
\end{equation}
where $\nu$ is the kinematic viscosity and $\varepsilon(u)=\tfrac12(\nabla u+\nabla u^{T})$ is the rate of strain tensor.

\section*{3. Boundary conditions}
On the moving lid the velocity is prescribed,

\begin{equation}
u = u_{\mathrm{lid}} \quad \text{on } y=1 ,
\end{equation}
and the remaining walls, $x=0$, $x=1$ and $y=0$, are no-slip walls, $u=0$. Because the velocity is prescribed on the whole boundary, the pressure is determined only up to an additive constant.

\section*{4. Parameters}
The viscosity is $\nu=0.1$ and the lid moves in the positive $x$-direction with $u_{\mathrm{lid}}=(1,0)$.

\section*{5. Discussion}
A primary vortex forms slightly above the centre of the cavity; secondary vortices appear in the lower corners only for much larger cavities.

\end{document}